\section{多重参考与区间对比}
本节介绍多条参考线、子弹图和多组区间对比。所有示例数据仅用于演示；参考标记和背景不改变原始数值及条形尺度。

\subsection{多条参考线}
\key{references} 接受一组记录，每条记录按
\key{数值/颜色/线型/名称} 填写，用逗号分隔并分别加花括号。
线型可选 \key{solid}、\key{dashed}、\key{dotted}。
数值必须在 min/max 内；名称可留空，但最后一个斜杠必须保留。
名称交替采用两层高度，长名称或密集参考仍需增加 width，或在表头解释名称。

\begin{example}[sidebyside=false]{示例 25：基准、目标和警戒值}
\gradbar[width=85mm,height=9pt,precision=0,
  references={{40/gray/dotted/基准},
              {75/gradbarsTeal/dashed/目标},
              {95/gradbarsOrange/solid/警戒}}]{68}
\end{example}
\key{references={}} 清除继承的参考线；原有 \key{target} 仍可共存。
参考线适用于普通条、正负条、堆叠、双层条、哑铃及区间条，不用于趋势线。
参考名称不参与数值格式化；名称内不要使用作为字段分隔符的斜杠。

\subsection{子弹图：分级背景、实际值和目标}
\key{\gradbullet[选项]{实际值}} 绘制较高的背景和居中的细条。
必须提供 \key{bullet bands}，每条记录为 \key{上界/颜色}。
第一个区间从 min 开始，后续区间从前一个上界开始；上界必须严格递增，最后一个必须等于 max。
背景区间只是参考等级，不会自动改变实际值颜色，也不会自动判定优劣方向。

\begin{example}[sidebyside=false]{示例 26：完成度子弹图}
\gradbullet[width=90mm,max=120,precision=0,
  color=gradbarsBlue,
  bullet bands={{60/black!8},{90/black!16},{120/black!26}},
  references={{90/gradbarsTeal/dashed/目标},
              {110/gradbarsOrange/solid/挑战}}]{82}
\end{example}
默认总高为 12pt，实际条占高度的 40\%，上下居中。
外部数值保留原值；NA 表示缺失，零表示实际值为零。
可使用原有条件配色、目标线和误差标记。
背景采用显式颜色；黑白打印时请像上例一样选用不同灰度。
\key{bullet bands} 仅用于子弹图，背景等级建议在表头或正文中说明。

\clearpage
\subsection{多组区间对比}
\key{\gradintervals[选项]{记录列表}} 将多组区间排列在同一个单元格内。
每条记录按 \key{名称/下界/点估计/上界/颜色} 填写。
所有行共用 width、min、max；左侧名称自动按最长名称对齐。
点估计为 NA 时不画圆点，区间仍保留；上下界同时为 NA 表示整组缺失。

\begin{example}[sidebyside=false]{示例 27：同一指标的三组不确定区间}
\gradintervals[width=72mm,min=0,max=100,
  height=5pt,precision=0,row padding=4pt,
  references={{70/gray/dashed/}}]{
  {方法 A/45/58/72/gradbarsBlue},
  {方法 B/55/68/82/gradbarsTeal},
  {方法 C/62/NA/88/gradbarsOrange}}
\end{example}
圆点为显式提供的点估计，外部标签为原始上下界；不从上下界推断均值或置信水平。
若区间表示 95\% 置信区间，应在表头或图注中写明。
点估计必须位于原始区间内，下界不得大于上界。
越界处理继承 \key{overflow}，所有组必须采用相同单位。

\begin{example}[sidebyside=false]{示例 28：正负区间与缺失组}
\begin{tabular}{ll}
\multicolumn{1}{c}{指标} & \multicolumn{1}{c}{估计区间}\\
变化量 & \gradintervals[min=-40,max=60,width=65mm,
  height=4pt,precision=0,row padding=3pt]{
  {组 A/-25/-10/5/gradbarsBlue},
  {组 B/5/18/35/gradbarsTeal},
  {组 C/NA/NA/NA/gray}}\\
\end{tabular}
\end{example}
该命令使用内部小表格，整体作为一个单元格排版，不在组间自动分页。
名称列、图区及外部区间标签都占用宽度；长名称请缩写，或调小图区 width。
参考名称会在每组重复，紧凑表格推荐留空，并在共同表头解释参考线。
记录名称及颜色字段不应包含斜杠；带逗号的名称应加额外花括号保护。
